\chapter*{幂指对专题练习}
\begin{enumerate}
    \item 幂函数的性质
    \begin{enumerate}
        \item 如图曲线是幂函数$y=x^n$在第一象限内的图像, 已知$n$取$\pm2,\pm\dfrac{1}{2}$四个值, 则曲线$c_1,c_2,c_3,c_4$对应的值依次为\dotfill(\kh{B})
        \xx{$-2,-\dfrac{1}{2},\dfrac{1}{2},2$}
        {$2,\dfrac{1}{2},-\dfrac{1}{2},-2$}
        {$-\dfrac{1}{2},-2,2,\dfrac{1}{2}$}
        {$2,\dfrac{1}{2},-2,-\dfrac{1}{2}$}
        \begin{minipage}{\linewidth}
        \centering
        
\includegraphics[totalheight=3cm]{./bodyxb/src/main42chap21mihanshu01.png}
        \end{minipage}
        \item 如图是函数$y=x^{\frac{m}{n}}$($m,n\zn$且互质)的图像, 则\dotfill(\kh{C})
        \xx{$m,n$是奇数, 且$\dfrac{m}{n}<1$}
        {$m$是偶数, $n$是奇数, 且$\dfrac{m}{n}>1$}
        {$m$是偶数, $n$是奇数, 且$\dfrac{m}{n}<1$}
        {$m$是奇数, $n$是偶数, 且$\dfrac{m}{n}>1$}
        \begin{minipage}{\linewidth}
        \centering
        \includegraphics[totalheight=3cm]{./bodyxb/src/main42chap21mihanshu02.png}
        \end{minipage}
        \item 下列命题中, 真命题的是\dotfill(\kh{C})
        \xx{幂函数中不存在既不是奇函数, 又不是偶函数的函数;}
        {如果一个幂函数不是偶函数, 那么它一定是奇函数;}
        {图像不经过$(-1,1)$的幂函数, 一定不是偶函数;}
        {如果两个幂函数有三个公共点, 那么这两个函数一定相同.}
        \item 已知幂函数 $y=\left(m^2-m-1\right) x^{\frac{m}{3}}$ 的图像关于 $y$ 轴对称, 则 $m=$ \blkda{$2$}.
        \item 若幂函数$y = (m^2-3m+1)x^{-m^2+2m+3}$的定义域为$\mathbf{R}$，$m$的值为\blkda{0}
        \begin{jieda}
            $m^2-3m+1 = 1$解得$m = 0$或$3$，$m = 0$的时候，$y = x^3$，符合要求，当$m = 3$的时候，$y = x^0$，定义域不为$\mathbf{R}$.
        \end{jieda}
        \item 当$x \in (0, +\infty)$时，幂函数$f(x) = (m^2-m-1)x^{-5m-3}$为增函数，则实数$m =$\blkda{$-1$}
        \begin{jieda}
            由幂函数概念可知$m^2-m-1 = 1$，解得$m = -1 \mbox{或} 2$，又因为$f(x)$在$(0, +\infty)$为增函数，所以有$-5m-3 > 0$，因此$m = -1$
        \end{jieda}
        \item 已知$(a+1)^{-\frac{3}{5}}<(3-2a)^{-\frac{3}{5}}$, 求实数$a$的取值范围. 
        \begin{jieda}
        \begin{enumerate}
            \item 方法1 \\
            构建函数$f(x) = x^{-\frac{3}{5}}$，原不等式转化为$f(a+1) < f(3-2a)$ \\
            易知函数$f(x) = x^{-\frac{3}{5}}$有两个单调减区间，奇函数，在$x > 0$时$f(x) > 0$，在$x < 0$时$f(x) < 0$ 
            \begin{dingautolist}{172}
                \item 若$(a+1)(3-2a) > 0$，即$a \in (-1, \dfrac{3}{2})$，则$a+1$和$3-2a$位于同一个单调区间，此时根据单调性，$f(a+1) < f(3-2a)$被转化为$a+1 > 3-2a$。故$a \in \left(\dfrac{2}{3}, \dfrac{3}{2}\right)$
                \item 若$a+1< 0 < 3-2a$，即$a \in (-\infty, -1)$，则$a+1$和$3-2a$位于不同单调区间，并且函数的性质可知满足$f(a+1) < f(3-2a)$
            \end{dingautolist}
            综上，$a \in (-\infty, -1) \cup \left(\dfrac{2}{3}, \dfrac{3}{2}\right)$
            \item 方法2 \\
            构建函数$f(x) = x^{\frac{3}{5}}$，原不等式转化为$f\left(\dfrac{1}{a+1}\right) < f\left(\dfrac{1}{3-2a}\right)$ \\
            易知函数$f(x) = x^{\frac{3}{5}}$为$R$上单调递增的奇函数，因此$f\left(\dfrac{1}{a+1}\right) < f\left(\dfrac{1}{3-2a}\right)$即$\dfrac{1}{a+1} < \dfrac{1}{3-2a}$ \\
            移项通分有$\dfrac{2-3a}{(a+1)(3-2a)} < 0$，解集为$a \in (-\infty, -1) \cup \left(\dfrac{2}{3}, \dfrac{3}{2}\right)$
        \end{enumerate}
        \end{jieda}
        \item 已知$m$是正整数，幂函数$f(x) = x^{1+m-m^2}$在$(0, +\infty)$上时严格增函数.
        \begin{enumerate}
            \item 求幂函数$f(x)$的解析式
            \item 设$\varphi(x) = \dfrac{a[f(x)]^2+1}{bf(x)+c}$（$a, b, c$为整数），若$\varphi(x)$是奇函数，又$\varphi(1) = 2, \varphi(2) < 3$，求$a, b, c$的值.
        \end{enumerate}
        \begin{jieda}
            \begin{enumerate*}
                \item 根据$f(x)$在$(0, +\infty)$上时严格增函数知$1+m-m^2>0$，又因为$m$是正整数，所以$m = 1$,故$f(x) = x$ \\
                \item $\varphi(x) = \dfrac{ax^2+1}{bx+c}$. 因为$\varphi(x)$是奇函数，故$0$不在定义域内，否则$\dfrac{1}{c}$必不为0，因此$c = 0$，经检验，此时$\varphi(x)$为奇函数. $\varphi(1) = \dfrac{a+1}{b} = 2$，即$b = \dfrac{a+1}{2}$；$\varphi(2) = \dfrac{4a+1}{2b} = \dfrac{4a+1}{a+1} = 4 - \dfrac{3}{a+1} < 3$，故$a \in \{0, 1\}$，又因为$b$是整数，故$a = 1$\\
                综上，$a = 1, b = 1, c = 0$
            \end{enumerate*}
        \end{jieda} 
    \end{enumerate}
    \item 对数运算 
    \begin{enumerate}
        \item 若$\log_427=m$, $\log_325=n$, 则$\lg2$可用$m,n$表示为\blkda{$\dfrac{3}{3+mn}$}. 
        \begin{jieda}
            $m = \dfrac{3log_23}{2}$，$n = \dfrac{2log_25}{log_23}$，$lg2 = \dfrac{1}{1+log_25}$，故$log_25 = \dfrac{m\cdot n}{2}$，$lg2 = \dfrac{3}{3+mn}$ \\
            \ps 对数运算表示类的题目，首先要统一底数，选取题干中的一个质因子作为底数，然后将所有的对数式都拆彻底（底数和真数互质），然后再凑.
        \end{jieda}
        \item 求$\dfrac{1-2log_65 \cdot log_{10}3 \cdot log_{15}2}{log_65 \cdot log_{10}3 + log_{10}3 \cdot log_{15}2 + log_{15}2 \cdot log_65}$的值.
        \begin{jieda}
            考虑换底，$log_65 = \dfrac{log_25}{log_23+1}, log_{10}3 = \dfrac{log_23}{log_25+1}, log_{15}2 = \dfrac{1}{log_23 + log_25}$\\
            故设$a = log_23, b = log_25$\\
            原式等于$\dfrac{1 - 2 \cdot \dfrac{b}{a+1} \cdot \dfrac{a}{b+1} \cdot \dfrac{1}{a+b}}{\dfrac{b}{a+1}\cdot \dfrac{a}{b+1} + \dfrac{a}{b+1} \cdot \dfrac{1}{a+b} + \dfrac{1}{a+b} \cdot \dfrac{b}{a+1}} = \dfrac{(a+1)(b+1)(a+b) - 2ab}{ab(a+b) + a(a+1) + b(b+1)} = 1$
        \end{jieda}
        \item 已知集合$U=\bold{R}$, $A = \{x, xy, lg(xy)\}$, $B = \{0, |x|, y\}$，若集合$A \cap B = B$，则$\left(x+\dfrac{1}{y}\right) + \left(x^2+\dfrac{1}{y^2}\right) + ... + \left(x^{2023}+\dfrac{1}{y^{2023}}\right) = $\blkda{-2}
        \begin{jieda}
            由$A \cap B = B$可知$B \subseteq A$，又由于集合中元素的互异性可知$A=B$.\\
            从$A$中含有元素$lg(xy)$可知，$x \neq 0$ 且 $y \neq 0$。由于$A=B$，可知 $A$中必然含有元素$0$。因此$lg(xy) = 0$, 即$xy = 1$。\\
            因此集合$B$中含有元素$1$，若有$x = |x|, y = 1$，会使得集合$A$违背互异性，因此必有$x < 0, y = |x|$，可知$x=y=-1$。即有$A = B = \{-1, 0, 1\}$。\\
            因此$\mbox{原式} = (-2)+2+(-2)+2+...+(-2) = -2$
        \end{jieda}
        \item 若正数$a, b$满足$log_2a + log_4b = 8, log_4a + log_8b = 2$，则$log_8a + log_2b = $ \blkda{$-\dfrac{52}{3}$}
        \begin{jieda}
            首先考虑统一底数为2，因此可设$A = log_2a, B = log_2b$。则有$log_4a = \dfrac{A}{log_24} = \dfrac{A}{2}$，类似可以推导其他的对数式。因此有
            $\left \{
                \begin{array}{cc}
                     A + \dfrac{B}{2} = 8  \\
                     \dfrac{A}{2} + \dfrac{B}{3} = 2
                \end{array}
                \right.$，求$\dfrac{A}{3} + B$。解二元一次方程得$A = 20, B = -24$，$\dfrac{A}{3} + B = -\dfrac{52}{3}$
        \end{jieda}
        \item 已知$log_2(x+y) = log_2x+log_2y$ \\
        \begin{enumerate*}[(1)]
            \item 求$x+2y$的最小值
            \item 求$\dfrac{4x}{x-1} + \dfrac{9y}{y-1}$的最小值
            \item 正数$z$满足$\dfrac{1}{x+y} + \dfrac{1}{z} = 1$，求$z$的取值范围
        \end{enumerate*}
        \begin{jieda}
            $log_2(x+y) = log_2(xy)$，故$x+y = xy$，即$x = \dfrac{y}{y-1}$，$y = \dfrac{x}{x-1}$，其中$x > 1, y > 1$.
            \begin{enumerate}[(1)]
                \item $x + 2y = x + \dfrac{2x}{x-1} = (x-1) + \dfrac{2}{x-1}+3 \geq 2\sqrt{2} + 3$，在$x = 1+\sqrt{2}, y = 1 + \dfrac{\sqrt{2}}{2}$时取等号.
                \item $\dfrac{4x}{x-1} + \dfrac{9y}{y-1} = \dfrac{4x}{x-1} + 9x = 9(x-1) + \dfrac{4}{x-1}+13 \geq 25$，在$x = \dfrac{5}{3}, y = \dfrac{5}{2}$时取等号.
                \item $\dfrac{1}{z} = 1 - \dfrac{1}{x+y} = 1 - \dfrac{x-1}{x^2} = 1 - \dfrac{1}{x} + \dfrac{1}{x^2}$\\
                令$t = \dfrac{1}{x}, t \in (0, 1)$，因此$\dfrac{1}{z} \in [\dfrac{3}{4}, 1)$，故$z \in (1, \dfrac{4}{3}]$
            \end{enumerate}
        \end{jieda}
    \end{enumerate}
    \item 对数函数的定义域
    \begin{enumerate}
        \item 已知集合$A = \{x | y = ln(x+1)\}$, $B = \{x | x^2 - 4 \leq 0\}$， 则$A \cup B = $\dotfill{(\kh{B})}. 
        \xx{$\{x | x \geq -2\}$}{$\{x | -1 < x \leq 2\}$}{$\{x | -1 < x < 2\}$}{$\{x | x \geq 2\}$}
        \begin{jieda}
            A是函数$y = ln(x+1)$的定义域，故$A=\{x|x>-1\}$，$B = \{x|-2 \leq x \leq 2\}$， 因此并集$A \cup B = \{x | x \geq -2\}$
        \end{jieda}
        \item 若$f(x) = (2x-1)^0 - \dfrac{log_2x}{\sqrt{x+2}}$，那么$f(x)$的定义域是\blkda[3]{$\left(0, \dfrac{1}{2}\right) \cup \left(\dfrac{1}{2}, +\infty\right)$}
        \begin{jieda}
            由题意可知，$2x-1 \neq 0, x > 0, x+2 > 0$， 故定义域为$(0, \dfrac{1}{2}) \cup (\dfrac{1}{2}, +\infty)$
        \end{jieda}
        \item 函数$lg\left(\dfrac{2-x}{x+3}\right)$的定义域为\blkda{$(-3, 2)$}
        \begin{jieda}
            根据对数函数的定义可知$\dfrac{2-x}{x+3} > 0$，解得$x \in (-3, 2)$
        \end{jieda}
        \item 函数$y=\sqrt{\log_{\frac{1}{3}}\left(\log_{\frac{1}{3}}\left(\log_{\frac{1}{3}}x\right)\right)}$的定义域是\blkda{$\left(\dfrac{1}{3},\sqrt[3]{\dfrac{1}{3}}\right]$}.
        \begin{jieda}
        $\log_{\frac{1}{3}}\left(\log_{\frac{1}{3}}\left(\log_{\frac{1}{3}}x\right)\right)\geq 0\iff 0<\log_{\frac{1}{3}}\left(\log_{\frac{1}{3}}x\right)\leq 1
        \iff 1>\log_{\frac{1}{3}}x\geq\dfrac{1}{3}\iff (\dfrac{1}{3})^1<x\leq(\dfrac{1}{3})^{\frac{1}{3}}$. 
        \end{jieda}
        \item  已知函数$y = lg(ax^2+ax+1)$，
        \begin{enumerate*}
            \item 若其定义域为$\mathbf{R}$，求参数$a$的取值范围
            \item 若其值域为$\mathbf{R}$，求参数$a$的取值范围
        \end{enumerate*}
        \begin{jieda}
            \begin{enumerate}
                \item 问题等价于$ax^2+ax+1 > 0$恒成立 \\
                $a=0$时显然成立; \\
                $a\neq 0$时, $a>0$且$ \Delta<0$, 即$a\in(0,4)$. \\
                故$a\in[0,4)$. 
                \item 问题等价于$y=ax^2+ax+1$的值域包含$(0,+\infty)$. \\
                $a=0$显然不符合条件; \\
                $a\neq 0$时, $a>0$且$\Delta\geq0$, 即$a\geq 4$. \\
                故$a\geq 4$.
            \end{enumerate}
        \end{jieda}
    \end{enumerate}
    \item 指数函数、对数函数的图像
    \begin{enumerate}
        \item 已知函数$f(x)=4+a^{x-1}$的图像恒过定点$P$, 则$P$点坐标是\blkda{$(1,5)$}. 
        \item 如图, 指数函数$y=a^x$, $y=b^x$, $y=c^x$, $y=d^x$, 则$a,b,c,d$与1的大小关系为\dotfill(\kh{B})
        \xx{$a<b<1<c<d$}
        {$b<a<1<d<c$}
        {$1<a<b<c<d$}
        {$a<b<1<d<c$}
        \begin{minipage}{\linewidth}
        \centering
        \includegraphics[totalheight=3cm]{./bodyxb/src/main42chap21zhishuhanshu01.png}
        \end{minipage}
        \ps 取横坐标为1，判断纵坐标即可.
        \item 函数$y=\log_2(x+1)$的图像不经过第\blkda{II、IV} 象限. 
        \begin{jieda}
        函数$y=\log_2(x+1)$的图像由$y=\log_2x$向左平移1个单位得到, 此时图像过原点, 作图易得不过第II、IV象限. 
        \end{jieda}
        
        \item 已知$1<x<2$, 则$a=2^x$, $b=\sqrt[3]{x}$, $c=\log_{\frac{1}{2}}x$的大小关系是\blkda{$a>b>c$}.
        \item 如果$\log_a2>\log_b2$, 那么$a,b$间的关系是\dotfill(\kh{D}) 
        \xx{$b>a>1$}
        {$1>b>a>0$}
        {$a>1$且$0<b<1$}
        {以上三种答案均有可能}
        \begin{jieda}
        \begin{enumerate}
            \item 方法1 \\
            通过换底公式可以转化为$\dfrac{\lg2}{\lg{a}}>\dfrac{\lg2}{\lg{b}}$，即$\dfrac{1}{\lg{a}}>\dfrac{1}{\lg{b}}$，因此三种可能均可满足要求.
            \item 方法2 \\
            结合图像可以判断出三种可能均满足要求 \\
            \ps 对于指数函数$y = a^x$，当$a > 1$时，$a$越大，函数图像增长越快，越贴近$y$轴；当$1 > a > 0$时，$\dfrac{1}{a}$越大，函数图像下降越快，越贴近$y$轴；\\
            对于对数函数$y = log_ax$，当$a > 1$时，$a$越大，函数图像增长越慢，越贴近$x$轴；当$1 > a > 0$时，$\dfrac{1}{a}$越大，函数图像下降越慢，越贴近$y$轴；
        \end{enumerate}
        \end{jieda}
    \end{enumerate}
    \item 指数函数和幂函数的交点
    \begin{enumerate}
        \item 函数$y = x^2$与$y = 2^x$的交点个数为\blkda{$3$}. 
        \begin{jieda}
            指数函数增长得很快（$a > 1$），当其从上方穿越了一个多项式函数后一定会从下方再次穿越这个多项式函数，$0 < a < 1$则是对称的情况.
        \end{jieda}
        \item 设常数$a \in \mathbf{R}$，则方程$|x+a| \cdot 2^x = 2025$的解的个数组成的集合是\blkda{$\{1,2,3\}$}
        \begin{jieda}
            方程变形为$|x+a| = 2025 \cdot 2^{-x}$，平移图像可知交点个数可能为1,2,3三种.
        \end{jieda}
        \item 设$min\{a, b, c\}$为$a, b, c$中最小值，则$y = min\{x+2, 10-x, 2^x\}$的最大值为\blkda{$6$}
        \begin{jieda}
            计算可知$y = 2^x$与$y = x+2$右侧交点为$(2, 4)$，故当$x > 2$时$x+2 < 2^x$，而$y = 10-x$与$y = 2-x$交于$(4, 6)$，结合图像可知$y = min\{x+2, 10-x, 2^x\}$的最大值为6.
        \end{jieda}
    \end{enumerate}
    \item 指数型复合函数的奇偶性
    \begin{enumerate}
        \item 已知函数$f(x)=\dfrac{a\cdot2^x+a^2-2}{2^x-1}(x\neq0)$, 其中$a$为实数. 
        \begin{enumerate}[(1)]
            \item 若函数$f(x)$为奇函数, 求出$a$的值;
            \item 若函数$f(x)$在$(0,+\infty )$上是增函数, 求$a$的取值范围;
            \item 当$x\in\left[1,2\right]$时, $f(x)>1$恒成立, 求$a$的取值范围. 
            \end{enumerate}
            \begin{jieda}
            $f(x)=a+\dfrac{a^2+a-2}{2^x-1}$
            \begin{enumerate}
            \item 根据$f(1)+f(-1)=0$, 可得$a=-1$或2. 经检验两者均满足要求：\\
            当$a=-1$时, $f(x)=-\dfrac{2^x+1}{2^x-1}$; 当$a=2$时, $f(x)=2\cdot\dfrac{2^x+1}{2^x-1}$.
            \item 设$g(x)=\dfrac{1}{2^x-1}$，则$f(x) = (a^2+a-2) \cdot g(x)+a$，由于$g(x)$单调减，故$a^2+a-2\leq 0$, 即$a\in[-2,1]$，可以使得$f(x)$单调增. 
            
            \item $f(x)$必然是单调函数，因此只需其两个端点满足要求即可保证恒成立，即$f(1) > 1, f(2) > 1$, 解得$a\in(-\infty,-5)\cup(1,+\infty)$. 
            \end{enumerate}
        \end{jieda}
        \item 若函数$f(x)=\dfrac{a-\me^x}{1+a\cdot\me^x}$($a$为常数)在定义域上为奇函数, 则$a$的值为\blkda{$\pm 1$}. 
        \begin{jieda}
        首选解法：$f(0)=0$或无意义，即$\dfrac{a-1}{1+a} = 0$或$\dfrac{a-1}{1+a}$无意义，故$a = \pm 1$\\
        也可以采用就极限思想：让函数在$\pm \infty$时的极限互为相反数求解：$x \rightarrow +\infty, f(x) \rightarrow -\dfrac{1}{a}, x \rightarrow -\infty, f(x) \rightarrow a$，由$a = \dfrac{1}{a}$得$a = \pm 1$. \\
        把握定义求解有：$f(x) = \dfrac{a-\me^x}{1+a\cdot\me^x}, f(-x) = \dfrac{a-\me^{-x}}{1+a\cdot\me^{-x}} = \dfrac{a\me^{x}-1}{a+\me^{x}}$ \\
        由$f(x) = -f(-x)$知$\dfrac{a-\me^x}{1+a\cdot\me^x} = \dfrac{1-a\me^{x}}{a+\me^{x}}$，交叉相乘得到$a^2 - e^{2x} = 1-a^2e^{2x}$，移项整理得$(a^2-1) = (1-a^2)\cdot(e^{2x})$，因为此式恒成立，所以必有$a = \pm 1$
        \end{jieda}
        \item 已知函数$f(x) = 2^x - (\frac{1}{2})^x+3$，且$f(m^2)+f(-5m+4)<6$，则实数$m$的取值范围是\blkda{$(1, 4)$}.
        \begin{jieda}
            可以设$g(x) = f(x) - 3$，将$f(m^2)+f(-5m+4)<6$转化为$g(m^2)+g(-5m+4)<0$后按照奇函数性质处理.
        \end{jieda}
            
        \item 已知$\mathbf{R}$上的奇函数$f(x)$与偶函数$g(x)$满足$f(x) - g(x) = 2^x$，若对于$\forall x \in (0, 2]$，$mf(x) - g(2x) \geq 0$恒成立，则$m$的取值范围是\blkda[2]{$[-2\sqrt{2}, +\infty)$}
        \begin{jieda}
            $f(x) = \frac{2^x - 2^{-x}}{2}, g(x) = -\frac{2^x+2^{-x}}{2}$ \\
            故$m(\frac{2^x - 2^{-x}}{2}) + \frac{4^x+4^{-x}}{2} \geq 0$在$(0, 2]$上恒成立. \\
            换元$t = 2^x - 2^{-x}, t \in (0, \frac{15}{4}]$，则$4^x+4^{-x} = t^2+2$，故有$mt + t^2 + 2 \geq 0$在$(0, \frac{15}{4}]$上恒成立. \\
            参变分离有$m \geq -(t + \frac{2}{t})$在$(0, \frac{15}{4}]$上恒成立.由对勾函数性质可知$-(t + \frac{2}{t})$在$(0, \frac{15}{4}]$上最大值为$-2\sqrt{2}$，因此有$m \in [-2\sqrt{2}, +\infty)$ \\
            \ps \textbf{一个定义域关于原点对称的函数$f(x)$一定能够拆成一个奇函数和一个偶函数之和：$f(x) = \frac{f(x) - f(-x)}{2} + \frac{f(x)+f(-x)}{2}$}
        \end{jieda}
    \end{enumerate}
    \item 指数型、对数型复合函数的单调区间
    \begin{enumerate}
        \item 函数$y=\left(\dfrac{1}{3}\right)^{-2x^2-8x+1}$单调递增区间是\blkda{$[-2,+\infty)$}. 
        \item 函数$f(x) = \left(\dfrac{1}{2}\right)^{\sqrt{x^2-3x-10}}$的递增区间是\blkda[3]{$(-\infty, -2]$}. 
        \begin{jieda}
            \textbf{定义域先行}：$x \in (-\infty, -2] \cup [5, +\infty)$ \\
            外层函数为单调减的指数函数，故要寻找整体的严格增区间，应寻找内层函数的严格减区间即$(-\infty,-2]$
        \end{jieda}
        \item 若函数$y=\dfrac{a^x}{2a^2-5a+2}$在$\bR$上为严格减函数, 则$a$的取值范围是\blkda[3]{$\left(0,\dfrac{1}{2}\right)\cup \left(1,2\right)$}. 
        \begin{jieda}
        $\begin{cases}
        a\in(0,1)\\
        2a^2-5a+2>0
        \end{cases}$
        或$\begin{cases}
        a>1\\
        2a^2-5a+2<0
        \end{cases}$, 
        解得$a\in\left(0,\dfrac{1}{2}\right)\cup(1,2)$. 
        \end{jieda}
        \item 设函数$f(x)=\dfrac{a(a^{2x}-1)}{a^x(a^2-1)}$($a>0$且$a\neq 1$), 证明: $y=f(x)$是$\bR$上为严格增函数. 
        \begin{jieda}
        $f(x)=\dfrac{a}{a^2-1}\cdot(a^x-a^{-x})$. \\
        任取实数集内的$x_1, x_2, x_1 < x_2$ \\
        $f(x_2)-f(x_1)=\dfrac{a}{a^2-1}\left(a^{x_2}-a^{x_1} - (a^{-x_2} - a^{-x_1})\right)$. 
        \begin{enumerate}
            \item 当$0 < a < 1$时,$a^{x_2}-a^{x_1} < 0$，$a^{-x_2} - a^{-x_1} > 0$，$\dfrac{a}{a^2-1} < 0$，故$f(x_2)-f(x_1) > 0$
            \item 当$a > 1$时,$a^{x_2}-a^{x_1} > 0$，$a^{-x_2} - a^{-x_1} < 0$，$\dfrac{a}{a^2-1} > 0$，故$f(x_2)-f(x_1) > 0$
        \end{enumerate}
        因此$f(x)$是$\bR$上为严格增函数. \\
        \ps $f(x)$并非恒正或者恒负，因此不应使用作商法.
        \end{jieda}
        \item 函数$f(x)=\log_2\left(x-\dfrac{a}{x}\right)(a>0)$在区间$[2,+\infty)$上是增函数, 则$a$的取值范围是\blkda{$(0,4)$}. 
        \begin{jieda}
        当$a>0$时, $y=x-\frac{a}{x}$在$(0,+\infty)$上单调递增. \\
        考虑到\textbf{定义域}, 于是$2-\frac{a}{2}>0$, 故$a\in(0,4)$. \\
        \ps 已知的单调区间必须是函数定义域的子集！
        \end{jieda}
        \item 已知函数$y = log_2(x^2 - ax - a)$在区间$(2, +\infty)$上严格增，则$a$的取值范围是\blkda{$\left(-\infty, \dfrac{4}{3}\right]$}. 
        \begin{jieda}
            外层函数严格增，故内层函数$y = x^2 - ax - a$满足在区间$(2, +\infty)$上严格增且函数值大于0.故应有$\dfrac{a}{2} \leq 2$且$4-3a \geq 0$，故$a \leq \dfrac{4}{3}$
        \end{jieda}
    \end{enumerate}
    
    \item 含指数、对数的分段函数严格单调
    \begin{enumerate}
        \item 已知函数$f(x) = \left \{ \begin{array}{ll}
         (2-5a)^x & x < 1  \\
         \dfrac{a}{x} + 5 & x \geq 1
        \end{array}\right.$ 在$\mathbf{R}$上严格递增，则实数$a$的取值范围是\blkda{$\left[-\dfrac{1}{2}, 0\right)$} .
        \begin{jieda}
            函数$f(x)$在$\mathbf{R}$上严格递增的充要条件为$\left \{ \begin{array}{l}
             2-5a > 1  \\
             a < 0 \\
             a + 5 \geq 2 - 5a
        \end{array}\right.$，解得$a \in \left[-\dfrac{1}{2}, 0\right)$
        \end{jieda}
        
        \item 已知函数 $y = \left\{  \begin{matrix}  - {x}^{2} + {2ax} - a,x \leq  1 \\  {2}^{x},x > 1 \end{matrix}\right.$ 在 $\mathbf{R}$ 上是严格增函数,则实数 $a$ 的取值范围是\blkda{$[1, 3]$}
        \begin{jieda}
            $\because y = \left\{  \begin{matrix}  - {x}^{2} + {2ax} - a,x \leq  1 \\  {2}^{x},x > 1 \end{matrix}\right.$ 在 $\mathbf{R}$ 上是严格增函数,
            
            $\therefore \left\{  \begin{matrix} a \geq  1 \\   - {1}^{2} + {2a} - a \leq  2 \end{matrix}\right.$ ,解得 $1 \leq  a \leq  3$ .
            
            故答案为: $\left\lbrack  {1,3}\right\rbrack$
        \end{jieda}

        \item 已知 $f\left( x\right)  = \left\{  \begin{matrix} \left( {2 - a}\right) x + 1,\left( {x < 1}\right) \\  {a}^{x},\left( {x \geq  1}\right)  \end{matrix}\right.$ 是 $\mathbf{R}$ 上的严格增函数,那么实数 $a$ 的取值范围是\blkda{$\left( {\dfrac{3}{2},2}\right)$}
        \begin{jieda}
            $f\left( x\right)$ 在 $\mathrm{R}$ 上单调递增,所以 $\left\{  {\begin{matrix} 2 - a > 0 \\  a > 1 \\  f\left( 1\right)  = a \geq  \left( {2 - a}\right)  \times  1 + 1 \end{matrix} \Longrightarrow  \left\{  {\begin{matrix} a < 2 \\  a > 1 \\  a \geq  \dfrac{3}{2} \end{matrix} \Longrightarrow  \dfrac{3}{2} \leq  a < 2}\right. }\right.$
        \end{jieda}
    \end{enumerate}
    \item 指数型复合函数的值域、最值
    \begin{enumerate}
    \item 函数$f(x)=\dfrac{3^x}{9^x+1}-\dfrac{1}{2}$的值域是\blkda{$\left(-\dfrac{1}{2},0\right]$}. 
    \begin{jieda}
    令$t = 3^x, t > 0, f(x) = g(t) = \dfrac{t}{t^2+1} - \dfrac{1}{2}$，转化为一次比二次求值域问题. \\
    $g(t) = \dfrac{t}{t^2+1} - \dfrac{1}{2} = \dfrac{1}{t+\dfrac{1}{t}} - \dfrac{1}{2} \in \left(-\dfrac{1}{2},0\right]$
    \end{jieda}
    \item 求函数$y=3\cdot4^x-20\cdot2^x+3\cdot4^{-x}-20\cdot2^{-x}$的最小值. 
    \begin{jieda}
    令$t=2^x+2^{-x}\in[2,+\infty)$, 则
    $y = 3t^2-20t-6 = 3(t-\dfrac{10}{3})^2-\dfrac{118}{3}$ \\
    故$y_{\min}=-\dfrac{118}{3}$, 当$t=\dfrac{10}{3}$即$2^x = \dfrac{1}{3}\mbox{或}3, x=log_23 \mbox{或}-log_23$时取到. 
    \end{jieda}
    \item 若$a^{2x}+\dfrac{1}{2}a^x-\dfrac{1}{2}\leq 0$($a>0$且$a\neq 1)$, 求函数$y=2a^{2x}-3a^x+4$的值域. 
    \begin{jieda}
    换元$t = a^x, t > 0$，原不等式化为$t^2+\dfrac{1}{2}t-\dfrac{1}{2}\leq 0$，解得$t\in(0,\dfrac{1}{2}]$. \\
    $y=2a^{2x}-3a^x+4=2(t-\dfrac{3}{4})^2+\dfrac{23}{8}\in[3,4)$. 
    \end{jieda}
    
    \item 设$F(x) = \dfrac{2^x}{a} - \dfrac{1}{2^x} + 2 - 2^{x-2} + \dfrac{a}{2^{x-2}}$，已知$F(x)$的最小值为$m$，且$m > 2 + \sqrt{7}$.求实数$a$的取值范围.
    \begin{jieda}
        $F(x) = \left(\dfrac{4-a}{4a}\right)\cdot2^x + \dfrac{(4a-1)}{2^x} + 2$，易知$a$会影响$F(x)$的单调性
        \begin{enumerate}
            \item $a \geq 4$ \\
            $\dfrac{4-a}{4a} \leq 0, (4a-1) > 0$，$F(x)$单调减，无最小值，不符合要求.
            \item $\dfrac{1}{4} < a < 4$ \\
            $\dfrac{4-a}{4a} > 0, (4a-1) > 0$，换元$t = 2^x (t > 0)$，则$F(x) = G(t) = \left(\dfrac{4-a}{4a}\right)t + \dfrac{(4a-1)}{t} + 2$为关于$t$的对勾函数的一支.因此$G(t) \geq 2\sqrt{\dfrac{(4-a)(4a-1)}{4a}}+2$，当$t = \sqrt{\dfrac{(4a-1)(4a)}{4-a}}$时取等号，\\
            故$2\sqrt{\dfrac{(4-a)(4a-1)}{4a}} > \sqrt{7}$，解得$a \in \left(\dfrac{1}{2}, 2\right)$ 
            \item $0 < a \leq \dfrac{1}{4}$ \\
            $\dfrac{4-a}{4a} > 0, (4a-1) \leq 0$，$F(x)$单调增，无最小值，不符合要求.
            \item $a = 0$\\
            $F(x)$无意义.
            \item $a < 0$ \\
            $\dfrac{4-a}{4a} < 0, (4a-1) < 0$，这种情况和第二种情况对称，此时只有最大值无最小值.
        \end{enumerate}
        综上，$a \in \left(\dfrac{1}{2}, 2\right)$
    \end{jieda}
    \end{enumerate}
    \item 对数型复合函数的值域、最值
    \begin{enumerate}
    \item 已知函数 $f\left( x\right)  = {\log }_{\frac{1}{2}}\left( {-{x}^{2} + {2x} + 7}\right)$ ,则 $f\left( x\right)$ 的值域是\blkda{$\lbrack  - 3, + \infty )$}
    \begin{jieda}
        令 $t =  - {x}^{2} + {2x} + 7,t > 0$ ,则 $y = {\log }_{\frac{1}{2}}t$ ,
        
        则 $t =  - \left( {{x}^{2} - {2x} + 1}\right)  + 8 =  - {\left( x - 1\right) }^{2} + 8 \leq  8$ ,可得 $0 < t \leq  8$ ,
        
        已知 $y = {\log }_{\frac{1}{2}}t$ 单调递减,所以 $y \geq  {\log }_{\frac{1}{2}}8 =  - 3$ ,
        
        则 $f\left( x\right)$ 的值域为 $\lbrack  - 3, + \infty )$ .
    \end{jieda}
    \item 已知 $f\left( x\right)  = 1 + {\log }_{3}x\left( {1 \leq  x \leq  9}\right)$ ,设 $g\left( x\right)  = {f}^{2}\left( x\right)  + f\left( {x}^{2}\right)$ ,则函数 $y = g\left( x\right)$ 的值域为\blkda{$[2, 7]$}
    \begin{jieda}
        由题意得 $\left\{  \begin{matrix} 1 \leq  x \leq  9 \\  1 \leq  {x}^{2} \leq  9 \end{matrix}\right.$ ,则 $1 \leq  x \leq  3$ ,即 $g\left( x\right)  = {f}^{2}\left( x\right)  + f\left( {x}^{2}\right)$ 的定义域为 $\left\lbrack  {1,3}\right\rbrack$ \\
        故 $g\left( x\right)  = {f}^{2}\left( x\right)  + f\left( {x}^{2}\right)  = {\left( 1 + {\log }_{3}x\right) }^{2} + 1 + {\log }_{3}{x}^{2} = {\left( {\log }_{3}x\right) }^{2} + 4{\log }_{3}x + 2$ \\
        令 ${\log }_{3}x = t,\left( {t \in  \left\lbrack  {0,1}\right\rbrack  }\right)$ ,则 $y = {t}^{2} + {4t} + 2 = {\left( t + 2\right) }^{2} - 2$ ,函数 $y = {\left( t + 2\right) }^{2} - 2$ 在 $\left\lbrack  {0,1}\right\rbrack$ 上单调递增,故 $y \in  \left\lbrack  {2,7}\right\rbrack$ \\
        故函数 $y = g\left( x\right)$ 的值域为 $\left\lbrack  {2,7}\right\rbrack$
    \end{jieda}
    \item 若函数 $f\left( x\right)  = {\log }_{2}\left\lbrack  {k{x}^{2} + \left( {{2k} - 1}\right) x + \dfrac{1}{4}}\right\rbrack$ 的值域为 $\mathbf{R}$ ,则实数 $k$ 的取值范围为\blkda[3]{$\left\lbrack  {0,\dfrac{1}{4}}\right\rbrack   \cup  \lbrack 1, + \infty )$}
    \begin{jieda}
        将问题转化为 $k{x}^{2} + \left( {{2k} - 1}\right) x + \dfrac{1}{4}$ 能取尽所有的正数,然后再分 $k = 0$ 和 $k \neq  0$ 两种情况, 并结合函数的性质求解即可. \\
        
        $\because$ 函数 $f\left( x\right)  = {\log }_{2}\left\lbrack  {k{x}^{2} + \left( {{2k} - 1}\right) x + \dfrac{1}{4}}\right\rbrack$ 的值域为 $\mathbf{R}$ ,

        $\therefore g\left( x\right)  = k{x}^{2} + \left( {{2k} - 1}\right) x + \dfrac{1}{4}$ 能取尽所有的正数.

        \begin{dingautolist}{172}
            \item 当 $k = 0$ 时, $g\left( x\right)  =  - x + \dfrac{1}{4}$ ,能取尽所有的正数,符合题意；
            \item 当 $k \neq  0$ 时,要使 $g\left( x\right)  = k{x}^{2} + \left( {{2k} - 1}\right) x + \dfrac{1}{4}$ 能取尽所有的正数,
            
            则需满足 $\left\{  \begin{matrix} k > 0 \\  \Delta  = {\left( 2k - 1\right) }^{2} - k = 4{k}^{2} - {5k} + 1 \geq  0 \end{matrix}\right.$ ,解得 $0 < k \leq  \dfrac{1}{4}$ 或 $k \geq  1$
        \end{dingautolist}
        
        综上可得 $0 \leq  k \leq  \dfrac{1}{4}$ 或 $k \geq  1$ ,
        
        $\therefore$ 实数 $k$ 的取值范围为 $\left\lbrack  {0,\dfrac{1}{4}}\right\rbrack   \cup  \lbrack 1, + \infty )$
    \end{jieda}
    \item 已知$x$满足不等式$2(\log_{\frac{1}{2}}x)^2+7\log_{\frac{1}{2}}x+3\leq 0$, 
    求$f(x)=\log_2\left(\dfrac{x}{2}\right)\log_2\left(\dfrac{x}{4}\right)$的最大值与最小值.
    \begin{jieda}
    令$t=\log_2x$, 则$2t^2-7t+3\leq 0$, 解得$t\in[\dfrac{1}{2},3]$. \\
    故$y=(t-1)(t-2)=t^2-3t+2=(t-\dfrac{3}{2})^2-\dfrac{1}{4}\in[-\dfrac{1}{4},2]$. \\
    所以$f(x)_{\min}=f(2^{\frac{3}{2}})=-\dfrac{1}{4}$, $f(x)_{\max}=f(8)=2$. 
    \end{jieda}
\end{enumerate}
    \item 指数函数、对数函数翻折
    \begin{enumerate}
        \item 已知定义在$R^{+}$上的函数$f(x) = \left \{\begin{array}{ll}
        |log_3x - 1|, & 0 < x \leq 9 \\
        4-\sqrt{x}, & x > 9
        \end{array} \right.$，设$a, b, c$是三个互不相同的实数，满足$f(a) = f(b) = f(c)$，求$abc$的取值范围.
        \begin{jieda}
            不妨设$a < b < c$，易知$0 < a < 3 < b < 9< c$，因此有$1 - log_3a = log_3b - 1$，即$ab = 9$. \\
            故$f(a) \in (0, 1)$，因此$c \in (9, 16)$，因此有$abc \in (81, 144)$
        \end{jieda}
        \item 关于$x$的方程$|3^x-1| = 2a$只有一个解，则$a$的取值范围是\blkda[3]{$\left[\dfrac{1}{2}, +\infty\right) \cup \{0\}$}.
        \begin{jieda}
            画图即可，$2a = 1$时是函数的渐近线，可以取得. 不要忘记$a = 0$时候也是只有一个解！
        \end{jieda}
    \item 已知函数$f(x) = \left \{ \begin{array}{ll}
        |log_2x| & 0 < x < 2 \\
        x^2 - 8x + 13 & x \geq 2
    \end{array} \right.$，若$f(x) = a$有四个解$x_1, x_2, x_3, x_4$，则$x_1 + x_2 + x_3 + x_4$的取值范围是\blkda{$\left(10, \dfrac{21}{2}\right)$}.
    \begin{jieda}
        $x_1 \cdot x_2 = 1, x_3 + x_4 = 8$，$x_1 + x_2 + x_3 + x_4 = 8 + x_2 + \dfrac{1}{x_2}, x_2 \in (1, 2)$，故$x_1 + x_2 + x_3 + x_4 \in \left(10, \dfrac{21}{2}\right)$ \\
        \ps \textbf{对数函数翻折时，总有和常函数交点横坐标乘积为定值}.
    \end{jieda}
    \end{enumerate}
    \item 指数复合型函数结合对称性求值
    \begin{enumerate}
        \item 已知$f(x) = \dfrac{(2^x+1)^2}{2^x \cdot x}+1$在$[-2018, -1]\cup[1, 2018]$上最大值为$M$,最小值为$m$,则$M+m = $\blkda{$2$}.
        \begin{jieda}
        设$g(x) = \dfrac{(2^x+1)^2}{2^x \cdot x}$，则$g(x) = \dfrac{2^x+\dfrac{1}{2^x}+2}{x} = \dfrac{2^x+2^{-x}+2}{x}$，故$g(-x) = -g(x)$，即$g(x)$为奇函数，故$M+m = 2$ \\
        \ps \textbf{涉及最大值与最小值和的题目一般可以优先考虑对称性}.
        \end{jieda}
        \item 已知函数$f(x) = \dfrac{3^{x+1}+1}{3^x+1}$在$[-2021, 2021]$上的最大值和最小值分别为$M, m$，则$M+m = $\blkda{$4$}.
        \begin{jieda}
            $g(x) = f(x) - 2 = \dfrac{3^x - 1}{3^x+1}$为奇函数，故$f(x)$关于$(0, 2)$中心对称，$M+m = 4$ \\
            \ps 本题如果没有发现$f(x) - 2$是奇函数，利用单调性也是可以解决的，计算量相对大. \\
            \ps 其实$y=\dfrac{a^x-1}{a^x+1}$, $y=\dfrac{a^x+1}{a^x-1}$都是奇函数.
        \end{jieda}
    \end{enumerate}
    \item 指数不等式与对数不等式
    \begin{enumerate}
        \item 已知不等式$a^{x^2-2x+3} \leq a^6$，且$a > 1$，则实数$x$的取值范围是\blkda{$[-1, 3]$}
        \begin{jieda}
            根据指数函数的单调性可以将不等式转化为$x^2-2x+3 \leq 6$，求解可得$x \in [-1, 3]$
        \end{jieda}
        \item 已知集合$A = \{x \in N^* | -2 < ln(x+1) < 2\}$，则集合$A$的子集个数为\blkda{64}
        \begin{jieda}
            首先将集合A中不等式转化为$\dfrac{1}{e^2} < x+1 < e^2$，因此集合A = $\{1,2,3,4,5,6\}$，因此其子集个数为64个。
        \end{jieda}
        \item 不等式 $log_{\frac{1}{2}}x+2 > x-1$的解集为\blkda{$(0, 2)$}
        \begin{jieda}
            $y = log_{\frac{1}{2}}x+2$单调减，$y = x-1$单调增，两者交点为$(2, 1)$，结合对数函数定义域可知解集为$(0, 2)$.
        \end{jieda}
        \item  已知函数$f(x)=\lg(1+2^x+3^x+4^x+a\cdot5^x)$对一切$x\leq 2$有意义, 实数$a$的取值范围是\blkda{$a>-\dfrac{6}{5}$}
        \begin{jieda}
        $1+2^x+3^x+4^x+a\cdot5^x>0$恒成立, 即$a>-\dfrac{1+2^x+3^x+4^x}{5^x}$恒成立. \\
        而$y=-\dfrac{1+2^x+3^x+4^x}{5^x}$在$\bR$上单调递增, 故$a>-\dfrac{6}{5}$($x=2$时). 
        \end{jieda}
        \item 若函数$f(2^x-1)$的定义域为$[-1, 1]$，则函数$f(log_2x - 1)$的定义域为\blkda{$[\sqrt{2}, 4]$}
        \begin{jieda}
            因为函数$f(2^x-1)$的定义域为$[-1, 1]$，所以$f(x)$的定义域为$[-\dfrac{1}{2}, 1]$，故$f(log_2x - 1)$的定义域为$-\dfrac{1}{2} \leq log_2x - 1 \leq 1$的解集，即$[\sqrt{2}, 4]$
        \end{jieda}
    \end{enumerate}
    \item 对数函数与均值不等式结合问题
    \begin{enumerate}
        \item 已知函数$f(x) = log_3\left(\dfrac{x}{3}\right) \cdot log_3\left(\dfrac{x}{27}\right)$，若对不相等的实数$x_1, x_2$有$f(x_1) = f(x_2)$，求$\dfrac{1}{x_1} + \dfrac{9}{x_2}$的最小值为\blkda{$\dfrac{2}{3}$}
        \item 已知$log_2(x+y) = log_2x+log_2y$, $\dfrac{4x}{x-1} + \dfrac{9y}{y-1}$的最小值为\blkda{25}
        \begin{jieda}
            $log_2(x+y) = log_2(xy)$，故$x+y = xy$，即$x = \dfrac{y}{y-1}$，$y = \dfrac{x}{x-1}$，其中$x > 1, y > 1$.\\
            $\dfrac{4x}{x-1} + \dfrac{9y}{y-1} = \dfrac{4x}{x-1} + 9x = 9(x-1) + \dfrac{4}{x-1}+13 \geq 25$，在$x = \dfrac{5}{3}, y = \dfrac{5}{2}$时取等号.
        \end{jieda}
    \end{enumerate}
    \item 指对方程
    \begin{enumerate}
        \item 基础运算
        \begin{enumerate}
            \item $2^{3x} = \dfrac{1}{4}$
            \begin{jieda}
                原方程等价于$2^{3x} = 2^{-2}$，故$3x = -2$，因此有$x = -\dfrac{2}{3}$
            \end{jieda}
            \item $log_2(x-2) = log_2(x^2-x-5)$
            \begin{jieda}
                原方程等价于 $\left \{ \begin{array}{ll}
                     x-2 = x^2-x-5  \\
                     x - 2 > 0 \\
                     x^2-x-5 > 0
                \end{array}\right.$ 解得$x = 3$
            \end{jieda}
        \end{enumerate}
        \item 换元
        \begin{enumerate}
            \item $9^x-4 \cdot 3^x - 45 = 0$
            \begin{jieda}
                换元$t = 3^x, t > 0$，故$t^2-4t - 45 = 0$，故$t = 9$或$-5$（舍），故$x = 2$
            \end{jieda}
            \item $9^x + 9^{-x} - 4(3^x+3^{-x}) + 5 = 0$
            \begin{jieda}
                换元$t = 3^x+3^{-x}, t \geq 2$，故$9^x + 9^{-x} = t^2-2$，故$t^2-4t+3 = 0$，解得$t =3$或$1$（舍），故$3^x = \dfrac{3\pm\sqrt{5}}{2}$，即$x = log_3{\dfrac{3-\sqrt{5}}{2}}$或$log_3{\dfrac{3+\sqrt{5}}{2}}$
            \end{jieda}
            \item $(log_3x)^2+log_{\frac{1}{3}}(x^2)-3 = 0$
            \begin{jieda}
                $(log_3x)^2+log_{\frac{1}{3}}(x^2)-3 = 0$ \\
                首先统一底数，原方程等价于$(log_3x)^2-2log_3x-3 = 0$\\
                换元$t = log_3x, t \in \mathbf{R}$，故$t^2 - 2t - 3 = 0$，解得$t = -1$或$3$，因此$x = \dfrac{1}{3}$或$27$
            \end{jieda}
            \item $log_3(3^x-1) \cdot log_3(3^{x-1}-\dfrac{1}{3})=2$
            \begin{jieda}
                原方程等价于$log_3(3^x-1) \cdot log_3(\dfrac{1}{3}(3^{x}-1))=2$，即$log_3(3^x-1) \cdot (log_3(3^{x}-1)-1)=2$ \\
                换元$t = log_3(3^x-1), t \in \mathbf{R}$，故$t(t-1) = 2$，解得$t = 2$或$-1$ \\
                因此$3^x - 1 = 9$或$\dfrac{1}{3}$，故$x = log_310$或$log_3\dfrac{4}{3}$ 
            \end{jieda}
        \end{enumerate}
        \item 指对转化
        \begin{enumerate}
            \item $x^{lgx} = \dfrac{x^3}{100}$
            \begin{jieda}
                原方程等价于$lgx \cdot lgx = 3lgx - 2$ \\
                换元$t = lgx$，故$t^2 -3t+2 = 0$，解得$t = 1$或$2$，故$x = 10$或$100$
            \end{jieda}
            \item $3 \cdot  {4}^{x+ 2}  = 4 \cdot  {3}^{{x}^{2}}$
            \begin{jieda}
                原方程等价于${4}^{x+1}  = {3}^{{x}^{2}-1}$
    
                两侧同时取对数有$(x+1)log_34 = x^2 - 1$
    
                合并同类项有$(x+1)(x-1-log_34) = 0$
    
                故$x = -1$或$log_312$.
            \end{jieda}
        \end{enumerate}
        \item 换底
        \begin{enumerate}
            \item $9^x + 4^x = \dfrac{5}{2} \cdot 6^x$
            \begin{jieda}
                原方程等价于$(\dfrac{3}{2})^x + (\dfrac{2}{3})^x = \dfrac{5}{2}$，换元$t = (\dfrac{3}{2})^x, t > 0$，故$t + \dfrac{1}{t} = \dfrac{5}{2}$，故$t = 2$或$\dfrac{1}{2}$，故$x = log_{\frac{3}{2}}2$或$-log_{\frac{3}{2}}2$
            \end{jieda}
            \item $3^x + 4^x = 5^x$
            \begin{jieda}
                原方程等价于$(\dfrac{3}{5})^x + (\dfrac{4}{5})^x = 1$ \\
                因为$y = (\dfrac{3}{5})^x + (\dfrac{4}{5})^x $单调递减，因此原方程只有一根$x = 2$
            \end{jieda}
            \item $2 \cdot x^{lg2} \cdot 2^{lgx} - 7x^{lg2} - 4 = 0$
            \begin{jieda}
                首先统一底数有$2 \cdot x^{lg2} \cdot x^{lg2} - 7x^{lg2} - 4 = 0$ \\
                换元$t = x^{lg2} = 2^{lgx}, t > 0$，因此$2t^2-7t-4 = 0$，解得$t = 4$或$-\dfrac{1}{2}$（舍）\\
                故$x = 100$
            \end{jieda}
        \end{enumerate}
    \end{enumerate}
    \item 反函数的定点问题
    \begin{enumerate}
        \item 设函数$f(x)$存在反函数$f^{-1}(x)$，且函数$y = x - f(x)$的图像过点$(1, 2)$，则函数$y = f^{-1}(x) - x$的图像一定过点\blkda{$(-1, 2)$}
        \begin{jieda}
            $y = x - f(x)$，代入$x = 1, y = 2$，则$f(1) = -1$，故$f^{-1}(-1) = 1$. $x = -1$代入$y = f^{-1}(x) - x$，$y = 2$，故过$(-1, 2)$ \\
            \ps 反函数过定点问题本质是$f(x)$过$(a, b)$，则反函数过$(b, a)$，在此基础上，通过代入特值解决问题。
        \end{jieda}
        \item 已知 $f(x)=\dfrac{a-x}{x-a-1}$, 若函数 $y=f(x)$ 的反函数图像的对称中心是 $(-1,3)$, 则实数 $a$ 的值是\blkda{$2$}.
        \begin{jieda}
            由反函数的对称中心知$f(x)$的对称中心为$(3, -1)$，因此$a+1 = 3$，故$a = 2$
        \end{jieda}
    \end{enumerate}
    \item 反函数平移问题
    \begin{enumerate}
        \item 如果函数$y = f(x), x \in \mathbf{R}$的反函数为$y = f^{-1}(x), x \in \mathbf{R}$，则函数$y = f^{-1}(x-1)(x\in \mathbf{R})$的反函数为\blkda[3]{$y = f(x)+1$}.
        \begin{jieda}
            反函数和原函数一个向下平移，令一个向左平移相同长度（长度为负，方向取反），两者可保持关于$y = x$对称.
        \end{jieda}
        \item 已知 $f(x)=\sqrt{x-1}$, 函数 $y=f(x)$ 的反函数为 $y=f^{-1}(x)$, 若 $f^{-1}(x)-a=f(x+a)$有实数根, 则实数 $a$ 的取值范围为\blkda{$\left[\dfrac{3}{4}, +\infty\right)$}.
        \begin{jieda}
            $f^{-1}(x) = x^2 + 1 (x \geq 0)$. \\
            设$g(x) = f(x+a)$，则$g^{-1}(x) = f^{-1}(x)-a$，因为$g(x)$与$g^{-1}(x)$均为增函数，因此其交点必然在$y = x$上。由图象可知，$a$的取值范围形式应为$[m, +\infty)$，$m$为边界情况对应的值.边界情况为$x^2 + 1 - a = x$只有一解，此时$a = \dfrac{3}{4}$，故$a \in \left[\dfrac{3}{4}, +\infty\right)$ 
        \end{jieda}
    \end{enumerate}
\end{enumerate}